\subsection{Dimension of graded modules}

In this number, we denote by M a graded H-module (of type Z).

Then \(\mathbf{S}^{-1}\mathbf{M}\) is an \(\mathbf{S}^{-1}\mathbf{H}\)-module, and if we set \(\mathbf{M}_{\geqslant n} = \bigoplus_{i\geqslant n}\mathbf{M}_i\) , we see as in

no. 1 that the sequence of sets \(S^{-1}M_{\geqslant n}\) is an exhaustive and separated filtration on

\(S^{-1}M\) and that the canonical map \(M \to S^{-1}M\) induces an isomorphism from M onto the graded module \(\oplus S^{-1}M_{\geqslant n}/S^{-1}M_{\geqslant n+1}\) .

Lemma 3. --- Suppose H is generated by \(\mathbf{H}_{0} \cup \mathbf{H}_{1}\) and M generated by \(\bigoplus_{i \leqslant n_{0}} \mathbf{M}_{i}\) for a suitable integer \(n_{0}\). Then the filtration \((\mathbf{S}^{-1}\mathbf{M}_{\geqslant n})\) of \(\mathbf{S}^{-1}\mathbf{M}\) is good for the ideal \(\mathbf{S}^{-1}\mathbf{H}_{\geqslant 1}\) of \(\mathbf{S}^{-1}\mathbf{H}\) .

We have, for \(n \geqslant n_0\) , \(\mathbf{M}_{\geqslant n+1} = \mathbf{H}_1 \cdot \mathbf{M}_{\geqslant n}\) , hence \(\mathbf{S}^{-1} \mathbf{M}_{\geqslant n+1} = \mathbf{H}_1 \cdot \mathbf{S}^{-1} \mathbf{M}_{\geqslant n} = \mathbf{S}^{-1} \mathbf{H}_{\geqslant 1} \cdot \mathbf{S}^{-1} \mathbf{M}_{\geqslant n}\) .

PROPOSITION 4. --- Suppose H is Noetherian and M is finitely generated. Then \(\dim_{\mathbf{H}}(\mathbf{M}) = \dim_{\mathbf{S}^{-1}\mathbf{H}}(\mathbf{S}^{-1}\mathbf{M})\) .

Let a be the annihilator of the H-module M; it is a graded ideal of H. Since M is a finitely generated H-module, the annihilator of the S \(^{-1}\) H-module S \(^{-1}\) M is the ideal S \(^{-1}\) a of S \(^{-1}\) H. We have dim \(_{H}\) (M) = dim(H/a) and dim \(_{S^{-1}H}\) (S \(^{-1}\) M) = dim(S \(^{-1}\) H/S \(^{-1}\) a). Prop. 4 then follows from th. 1, c) of no. 2 applied to the graded ring H/a.

Proposition 5. — Suppose \(\mathbf{H}_{0}\) is local and Artinian, \(\mathbf{H}\) is generated by \(\mathbf{H}_{0} \cup \mathbf{H}_{1}\) , \(\mathbf{H}_{1}\) is of finite type as a \(\mathbf{H}_{0}\) -module, and \(\mathbf{M}\) is nonzero and of finite type as an \(\mathbf{H}\) -module. Then \(\mathbf{M}_{n}\) is a nonzero finitely generated \(\mathbf{H}_{0}\) -module of finite length for each \(n\) , and there exists \(\mathbf{Q}(\mathbf{T}) \in \mathbf{Z}[\mathbf{T}, \mathbf{T}^{-1}]\) such that \(\mathbf{Q}(1) > 0\) and that one has in the ring \(\mathbf{Z}((\mathbf{T}))\)

\[
\sum_ {n \in \mathbf {Z}} \mathrm{long} _ {\mathrm{H} _ {0}} (\mathrm{M} _ {n}). \mathrm{T} ^ {n} = (1 - \mathrm{T}) ^ {- d}. \mathrm{Q} (\mathrm{T}),
\]

with \(d = \dim_{\mathbf{H}}(\mathbf{M})\).

The ring \(\mathbf{S}^{-1}\mathbf{H}\) is local and Noetherian (no. 1, remark 1), the \(\mathbf{S}^{-1}\mathbf{H}\) -module \(\mathbf{S}^{-1}\mathbf{M}\) is nonzero of finite type, and of dimension \(d = \dim_{\mathbf{H}}(\mathbf{M})\) (prop. 4). Moreover, \(\mathbf{S}^{-1}\mathbf{H}_{\geqslant 1}\) is a defining ideal of \(\mathbf{S}^{-1}\mathbf{H}\) (§ 3, no. 2, Lemma 2) and \(\mathbf{S}^{-1}\mathbf{M}_{\geqslant n}\) is a filtration \(\mathbf{S}^{-1}\mathbf{H}_{\geqslant 1}\) -good on \(\mathbf{S}^{-1}\mathbf{M}\) (Lemma 3). Finally, we have long \(_{\mathbf{S}^{-1}\mathbf{H}}\mathbf{S}^{-1}\mathbf{M}_{\geqslant n}/\mathbf{S}^{-1}\mathbf{M}_{\geqslant n+1} = \text{long}_{\mathbf{H}_0}(\mathbf{M}_n)\) for every \(n\). It therefore suffices to apply theorems 2 and 3 of § 4 (nos. 3 and 4).

Remark. --- Except for the determination of the integer d, prop. 5 follows directly from th. 1 of § 4, no. 2.

COROLLARY. --- Let A be a Noetherian local ring and q an ideal of definition of A. Then we have \(\dim(A) = \dim(\text{gr}_{\mathfrak{q}}(A))\) .

Applying prop. 5 to the case M = H = gr\_q(A), we obtain the relation

\[
\sum_ {n \geqslant 0} \mathrm{long} _ {\mathrm{A} / \mathfrak {q}} (\mathfrak {q} ^ {n} / \mathfrak {q} ^ {n + 1}). \mathrm{T} ^ {n} = (1 - \mathrm{T}) ^ {- d} \mathrm{Q} (\mathrm{T})
\]

with \(d = \dim(\mathrm{gr}_{\mathfrak{q}}(\mathbf{A}))\) and \(\mathbf{Q}(1) \neq 0\) . We have \(d = \dim(\mathbf{A})\) by th. 3 of § 4, no. 4, whence the corollary.

PROPOSITION 6. --- Suppose \(\mathbf{H}_{0}\) is local and Artinian, \(\mathbf{H}\) is of finite type as a \(\mathbf{H}_{0}\) -algebra and \(\mathbf{M}\) is of finite type as a \(\mathbf{H}\) -module.

\begin{enumerate}
\def\labelenumi{\alph{enumi})}
\item
  Let \(a_1, \ldots, a_n\) be elements of H, homogeneous of degrees \(> 0\) , and let \(\varphi\) be the homomorphism (of \(\mathbf{H}_0\) -algebras) from \(\mathbf{H}_0[\mathbf{X}_1, \ldots, \mathbf{X}_n]\) to H which sends \(\mathbf{X}_i\) to \(a_i\) for \(1 \leqslant i \leqslant n\) . The \(S^{-1}H\) -module \(S^{-1}M / \sum_{i=1}^{n}(a_i/1) \cdot S^{-1}M\) is of finite length if and only if \(\varphi_*(M)\) is a finitely generated module over \(\mathbf{H}_0[\mathbf{X}_1, \ldots, \mathbf{X}_n]\) .
\item
  There exists a family \((a_{1},...,a_{d})\) of elements of H, all homogeneous of the same degree \textgreater0, with \(d=\dim_{\mathrm{H}}(\mathbf{M})\) and such that \((a_{1}/1,...,a_{d}/1)\) is a maximal secant sequence for the \(S^{-1}H\) -module \(S^{-1}M\) . If moreover H is generated by \(H_{1}\) as \(H_{0}\) -algebra, and if the residue field of \(H_{0}\) is infinite, one may take the \(a_{i}\) of degree 1.
\item
  Set \(\mathbf{N} = \mathbf{M} / \sum_{i=1}^{n} a_i \mathbf{M}\) . We have \(\dim_{\mathrm{H}}(\mathbf{N}) = \dim_{\mathrm{S}^{-1}\mathrm{H}}(\mathrm{S}^{-1}\mathrm{N})\) by prop. 4. Consequently, the \(\mathrm{S}^{-1}\mathrm{H}\) -module \(\mathrm{S}^{-1}\mathrm{N}\) is of finite length if and only if the H-module N is of finite length, that is, if and only if N is a \(\mathrm{H}_0\) -module of finite type. If \(\varphi_*(\mathbf{M})\) is the module over \(\mathrm{H}_0[\mathrm{X}_1, ..., \mathrm{X}_n]\) deduced from M by the homomorphism \(\varphi : \mathrm{H}_0[\mathrm{X}_1, ..., \mathrm{X}_n] \to \mathbf{H}\) , one has \(\mathbf{N} = \varphi_*(\mathbf{M}) / \sum_{i=1}^{n} X_i \cdot \varphi_*(\mathbf{M})\) . Consequently (A, II, p.~171, cor. 3 and remark) \(\varphi_*(\mathbf{M})\) is a finitely generated module over \(\mathrm{H}_0[\mathrm{X}_1, ..., \mathrm{X}_n]\) if and only if N is a \(\mathrm{H}_0\) -module of finite type. This proves a).
\end{enumerate}

To prove b), we will first establish a lemma.

Lemma 4. --- Let \(b\) be an element of H, homogeneous of degree \(>0\), and not belonging to any of the minimal elements \(\mathfrak{p}\) of \(\operatorname{Supp}(\mathbf{M})\) such that \(\dim(\mathrm{H}/\mathfrak{p}) = \dim_{\mathrm{H}}(\mathbf{M})\). We then have \(\dim_{\mathrm{H}}(\mathbf{M}/b\mathbf{M}) = \dim_{\mathrm{H}}(\mathbf{M}) - 1\).

Set \(d = \dim_{\mathbf{H}}(\mathbf{M})\). By the definition of \(b\) , one has \(\dim_{\mathbf{H}}(\mathbf{M}/b\mathbf{M}) < d\). By prop. 4, we have

\[
\dim_ {\mathrm{H}} (\mathrm{M} / b \mathrm{M}) = \dim_ {\mathrm{S} ^ {- 1} \mathrm{H}} \left(\mathrm{S} ^ {- 1} \mathrm{M} / (b / 1). \mathrm{S} ^ {- 1} \mathrm{M}\right)
\]

and formula (8) of § 3, no. 2 yields the inequality

\[
\dim_ {\mathrm{S} ^ {- 1} \mathrm{H}} (\mathrm{S} ^ {- 1} \mathrm{M} / (b / 1). \mathrm{S} ^ {- 1} \mathrm{M}) \geqslant \dim_ {\mathrm{S} ^ {- 1} \mathrm{H}} (\mathrm{S} ^ {- 1} \mathrm{M}) - 1.
\]

Finally, we have

\[
\dim_ {S ^ {- 1} H} (S ^ {- 1} M) = \dim_ {H} (M) = d
\]

by prop. 4. We therefore have \(\dim_{\mathrm{H}}(\mathrm{M}/b\mathrm{M}) \geqslant d - 1\) , whence lemma 4.

Let us resume the proof of prop. 6, b). We may assume \(\dim_{\mathbf{H}}(\mathbf{M}) > 0\). Note that every minimal element of \(\operatorname{Supp}(\mathbf{M})\) is graded (apply Lemma 1 of No. 2 to the quotient of H by the annihilator of M). By Prop. 8 of III, § 1, No. 4, there therefore exists a homogeneous element \(b\) of H, of degree \(>0\), not belonging to any of the minimal elements \(\mathfrak{p}\) of \(\operatorname{Supp}(\mathbf{M})\) such that \(\dim(\mathrm{H}/\mathfrak{p}) = \dim_{\mathbf{H}}(\mathbf{M})\). By Lemma 4, we have \(\dim_{\mathbf{H}}(\mathbf{M}/b\mathbf{M}) = \dim_{\mathbf{H}}(\mathbf{M}) - 1\). Suppose moreover that H is generated by \(\mathrm{H}_{1}\) as \(\mathrm{H}_{0}\) -algebra and the residue field \(k\) of \(\mathrm{H}_{0}\) is infinite. For every minimal element \(\mathfrak{p}\) of \(\operatorname{Supp}(\mathbf{M})\) such that \(\dim(\mathrm{H}/\mathfrak{p}) = \dim_{\mathbf{H}}(\mathbf{M})\) , consider the vector subspace \(\mathrm{V}_{\mathfrak{p}} = (\mathfrak{p} \cap \mathrm{H}_{1}) \otimes_{\mathrm{H}_{0}} k\) of the \(k\)-vector space \(\mathrm{V} = \mathrm{H}_{1} \otimes_{\mathrm{H}_{0}} k\). If one had \(\mathrm{V}_{\mathfrak{p}} = \mathrm{V}\) , one would have \(\mathfrak{p} \cap \mathrm{H}_1 = \mathrm{H}_1\) (II, §3, no. 2, prop. 4), whence \(\mathrm{H}_1 \subset \mathfrak{p}\) and \(\dim_{\mathrm{H}}(\mathbf{M}) = \dim(\mathrm{H}/\mathfrak{p}) \leqslant \dim(\mathrm{H}/\mathrm{H}_{\geqslant 1}) = 0\) , which is not the case. Since \(k\) is assumed to be infinite, the union of the \(\mathrm{V}_{\mathfrak{p}}\) is distinct from V; if \(b \in \mathrm{H}_1\) is such that \(b \otimes 1\) belongs to none of the \(\mathrm{V}_{\mathfrak{p}}\) , one has \(\dim(\mathrm{M}/b\mathrm{M}) = \dim_{\mathrm{H}}(\mathrm{M}) - 1\) .

Proceeding by induction on \(d = \dim_{\mathbf{H}}(\mathbf{M})\) , one then constructs a sequence \((b_1, \ldots, b_d)\) of elements of H, with \(b_i\) homogeneous of degree \(n_i > 0\) and such that \(\mathbf{M} / \sum_{i=1}^{n} b_i \mathbf{M}\) is an H-module of finite length. If one assumes H is generated by \(\mathbf{H}_1\) as \(\mathbf{H}_0\) -algebra and the residue field of \(\mathbf{H}_0\) infinite, one may assume \(n_i = 1\) for \(i = 1, \ldots, d\). By prop. 4, we have \(\dim_{\mathbf{S}^{-1}\mathbf{H}}(\mathbf{S}^{-1}\mathbf{M}) = d\) and

\[
\dim_ {\mathrm{S} ^ {- 1} \mathrm{H}} (\mathrm{S} ^ {- 1} \mathrm{M} / \sum_ {i = 1} ^ {d} (b _ {i} / 1). \mathrm{S} ^ {- 1} \mathrm{M}) = 0.
\]

Then \((b_{1} / 1,\dots,b_{d} / 1)\) is a maximal secant sequence for the \(\mathbf{S}^{-1}\mathbf{H}\) -module \(\mathbf{S}^{-1}\mathbf{M}\) . Put \(a_{i} = b_{i}^{(n_{1}\dots n_{d}) / n_{i}}\) for \(1\leqslant i\leqslant d\) ; then the \(a_{i}\) are all of the same degree, and \((a_{1} / 1,\dots,a_{d} / 1)\) is a maximal secant sequence for \(\mathbf{S}^{-1}\mathbf{M}\) (§ 3, no. 2, Remark 3).

COROLLARY 1. --- Suppose that \(\mathbf{H}_0\) is a field and that \(\mathbf{H}\) is of finite type as an \(\mathbf{H}_0\) -algebra. Set \(n = \dim(\mathbf{H})\) . There exist homogeneous elements \(a_1, ..., a_n\) of \(\mathbf{H}\) all of the same degree \(>0\) , such that the \(\mathbf{H}_0\) -homomorphism \(\varphi: \mathbf{H}_0[\mathbf{X}_1, ..., \mathbf{X}_n] \to \mathbf{H}\) defined by \(\varphi(\mathbf{X}_i) = a_i\) , \(i = 1, ..., n\) , be injective and make \(\mathbf{H}\) a finite \(\mathbf{H}_0[\mathbf{X}_1, ..., \mathbf{X}_n]\)-algebra. If \(\mathbf{H}\) is generated by \(\mathbf{H}_1\) as \(\mathbf{H}_0\) -algebra and if \(\mathbf{H}_0\) is infinite, one may assume the \(a_i\) of degree 1.

There exists, by prop. 6, a \(\mathbf{H}_0\) -homomorphism \(\varphi\) of the indicated form that makes H a \(\mathbf{H}_0[\mathbf{X}_1, \ldots, \mathbf{X}_n]\) -finite algebra. By th. 1 of § 2, no. 3, we then have

\[
\dim \left(\mathrm{H} _ {0} \left[ \mathrm{X} _ {1}, \dots , \mathrm{X} _ {n} \right] / (\operatorname{Ker} \varphi)\right) = \dim (\mathrm{H});
\]

since we have

\[
\dim (\mathrm{H}) = n = \dim \left(\mathrm{H} _ {0} \left[ \mathrm{X} _ {1}, \dots , \mathrm{X} _ {n} \right]\right),
\]

and that \(\mathrm{H}_0[\mathrm{X}_1, \dots, \mathrm{X}_n]\) is an integral domain, this implies \(\operatorname{Ker} \varphi = \{0\}\) .

Note. --- Let \((h_{1}, \ldots, h_{r})\) be a finite generating system of the \(H_{0}\) -vector space \(H_{1}\) . For \(\lambda = (\lambda_{1}, \ldots, \lambda_{r}) \in H_{0}^{r}\) , set \(h_{\lambda} = \lambda_{1} h_{1} + \cdots + \lambda_{r} h_{r}\). The proofs of Prop. 6 and Cor. 1 yield the following result: the set of elements \((\lambda_{1}, \ldots, \lambda_{n})\) of \((\mathrm{H}_{0}^{r})^{n}\) such that the elements \(a_{i} = h_{\lambda_{i}} \in H_{1}\) satisfy the conclusion of Cor. 1, contains the complement of the union of a finite number of vector subspaces of \((\mathrm{H}_{0}^{r})^{n}\) distinct from the entire space.

COROLLARY 2. --- Let A be a Noetherian local ring and \(n\in\mathbb{N}\). For one to have \(\dim(\mathbf{A})\geqslant n\) it is necessary and sufficient that for every integer \(r\geqslant0\) , one has

\[
\left[ \mathfrak {m} _ {\mathrm{A}} ^ {r} / \mathfrak {m} _ {\mathrm{A}} ^ {r + 1}: \kappa_ {\mathrm{A}} \right] \geqslant \binom {n + r - 1} {n - 1}, \quad \left(\text { resp.   long } _ {\mathrm{A}} (\mathrm{A} / \mathfrak {m} _ {\mathrm{A}} ^ {r + 1}) \geqslant \binom {n + r} {n}\right);
\]

We have equality for every r if and only if A is regular of dimension n.

The condition is sufficient (§ 4, no. 4, th. 3 and § 5, no. 2, th. 1). Let us show that it is necessary. Consider the graded ring gr(A) = gr \(_{m_A}\) (A); let k be an infinite extension of the field κ \(_A\) , and set H = k ⊗ κ \(_A\) gr(A). The ring H has dimension ≥ n (prop. 5 and its corollary); hence, from cor. 1, we deduce the existence of an injective graded homomorphism of graded k-algebras φ : H \(_0\) {[}X \(_1\) , \ldots, X \(_n\) {]} → H. Consequently, for every integer r ≥ 0,

\[
\left[ \mathrm{gr} _ {r} (\mathrm{A}): \kappa_ {\mathrm{A}} \right] = \left[ \mathrm{H} _ {r}: \mathrm{H} _ {0} \right] \geqslant \binom {n + r - 1} {n - 1},
\]

and the equality for every \(r\) implies the bijectivity of \(\varphi\) , hence the regularity of A ( \(\S\) 5, no. 2, th. 1).

The equalities

\[
\mathrm{long} _ {\mathbf {A}} (\mathrm{A} / \mathfrak {m} _ {\mathbf {A}} ^ {r + 1}) = \sum_ {i = 0} ^ {r} \left[ \mathrm{gr} _ {i} (\mathbf {A}): \kappa_ {\mathbf {A}} \right]
\]

and

\[
\binom {n + r} {n} = \sum_ {i = 0} ^ {r} \binom {n + i - 1} {n - 1}
\]

then imply the analogous assertions for the function \(r \mapsto \text{long}_{\mathbf{A}}(\mathbf{A}/m_{\mathbf{A}}^{r+1})\) .

